\chapter{Type 2 Diabetes}

\section*{Definition and Overview}
Type 2 diabetes (T2D) is a metabolic disorder characterised by hyperglycaemia resulting from progressive insulin resistance in liver, muscle, and adipose tissue, coupled with inadequate compensatory insulin secretion from pancreatic $\beta$-cells. It accounts for ~90\% of all diabetes. Onset is typically in adulthood (>40 years) but increasingly seen in adolescents and young adults due to rising obesity. T2D is a heterogeneous disease with varying contributions of insulin resistance and $\beta$-cell dysfunction. It is strongly associated with overweight/obesity, physical inactivity, and genetic susceptibility. Remission is achievable with substantial weight loss, particularly early in the disease course.

\section*{Epidemiology}
In Australia, ~1.3 million adults (5.3\%) have diagnosed T2D; an estimated 500,000 have undiagnosed T2D. Prevalence increases with age: 10--15\% at 65--74, ~20\% at $\ge$75. Indigenous Australians have 3--4x higher prevalence, earlier onset (median age 45 vs 65), and worse complications. Global prevalence has quadrupled since 1980 (>500 million). Risk factors: overweight/obesity (BMI $\ge$25: 3--7x risk), central adiposity, physical inactivity, family history, ethnicity (South Asian, Pacific Islander, Middle Eastern, Indigenous), gestational diabetes history, PCOS, antipsychotic use, low birth weight. Socioeconomic disadvantage doubles prevalence.

\section*{Aetiology and Pathophysiology}
\begin{itemize}
    \item \textbf{Insulin resistance:} reduced glucose uptake in muscle/adipose; failed suppression of hepatic glucose output; increased lipolysis $\to$ elevated FFAs $\to$ ectopic fat (liver, muscle, pancreas); inflammation (TNF-$\alpha$, IL-6, MCP-1) from adipose tissue macrophages
    \item \textbf{$\beta$-cell dysfunction:} progressive loss of $\beta$-cell mass and function; glucotoxicity, lipotoxicity, amyloid (IAPP), oxidative stress, ER stress; dedifferentiation/transdifferentiation; 50\% functional loss at diagnosis
    \item \textbf{Genetics:} >400 loci (TCF7L2 strongest); heritability ~30--50\%; polygenic risk scores; monogenic forms (MODY) distinct
    \item \textbf{Developmental origins:} fetal undernutrition $\to$ thrifty phenotype $\to$ mismatch with postnatal obesogenic environment; low birth weight + adult obesity = highest risk
    \item \textbf{Gut microbiome:} reduced diversity, altered metabolites (SCFAs, bile acids, TMAO); influences inflammation, GLP-1 secretion, energy harvest
    \item \textbf{Adipose tissue expandability hypothesis:} limited subcutaneous fat storage capacity $\to$ ectopic/visceral fat $\to$ metabolic syndrome
    \item \textbf{Natural history:} normal glucose $\to$ IFG/IGT (prediabetes) $\to$ T2D $\to$ $\beta$-cell failure $\to$ insulin requirement; trajectory varies widely
\end{itemize}

\section*{Clinical Features}
\begin{itemize}
    \item \textbf{Often asymptomatic:} detected on screening (HbA1c, fasting glucose) or incidentally
    \item \textbf{Symptomatic (when marked hyperglycaemia):} polyuria, polydipsia, fatigue, blurred vision, recurrent infections (candidiasis, UTIs, skin), slow wound healing, weight loss (if severe insulin deficiency)
    \item \textbf{Signs of insulin resistance:} acanthosis nigricans (axillae, neck), skin tags, central obesity, hypertension
    \item \textbf{Complications at diagnosis:} 20--30\% have retinopathy, 10--20\% nephropathy, 10--15\% neuropathy, 20--30\% CVD (reflects years of undiagnosed dysglycaemia)
    \item \textbf{Associated conditions:} NAFLD/MAFLD (70\%), OSA (50--70\%), PCOS, gout, erectile dysfunction, depression, cognitive impairment
\end{itemize}

\section*{Differential Diagnosis}
\begin{itemize}
    \item \textbf{Type 1 diabetes / LADA:} younger, leaner, autoantibody positive, rapid insulin requirement; test GAD/IA-2/ZnT8 if uncertain
    \item \textbf{MODY:} young (<25), lean, strong autosomal dominant family history, negative autoantibodies, measurable C-peptide; genetic testing
    \item \textbf{Secondary diabetes:} pancreatic (CF, chronic pancreatitis, haemochromatosis), drug-induced (steroids, antipsychotics, calcineurin inhibitors, checkpoint inhibitors), endocrine (Cushing, acromegaly, phaeochromocytoma)
    \item \textbf{Prediabetes (IFG/IGT):} fasting glucose 6.1--6.9 / 2-h OGTT 7.8--11.0; high risk for T2D; lifestyle intervention prevents progression
    \item \textbf{Stress hyperglycaemia:} acute illness; resolves; but predicts future T2D
    \item \textbf{Type 3c (pancreatogenic) diabetes:} exocrine pancreatic insufficiency, history of pancreatitis, low faecal elastase
\end{itemize}

\section*{When to Seek Medical Care}
Screening (HbA1c or fasting glucose) from age 40 (18 Indigenous) or earlier with risk factors (AUSDRISK $\ge$12). Urgent review for symptomatic hyperglycaemia, ketosis, or HHS (hyperosmolar hyperglycaemic state -- severe dehydration, marked hyperglycaemia $>$30 mmol/L, altered consciousness, no significant ketosis).

\textbf{Call 000 for:}
\begin{itemize}
    \item HHS: severe dehydration, altered consciousness, seizures
    \item DKA (uncommon in T2D but occurs with SGLT2i, severe illness, insulin deficiency)
    \item Severe hypoglycaemia (on insulin/sulfonylureas): unconsciousness, seizures
    \item Acute cardiovascular event (MI, stroke) -- hyperglycaemia common
\end{itemize}

\section*{Diagnosis and Investigations}
\begin{itemize}
    \item \textbf{Diagnostic criteria (any one, confirm with repeat unless symptomatic):} fasting glucose $\ge$7.0 mmol/L, random $\ge$11.1 with symptoms, HbA1c $\ge$6.5\% (48 mmol/mol), OGTT 2-h $\ge$11.1
    \item \textbf{HbA1c:} preferred for screening/diagnosis; no fasting required; reflects 3-month average; caveats: haemoglobinopathies, anaemia, CKD, pregnancy, recent blood loss/transfusion
    \item \textbf{Baseline assessment:} HbA1c, lipids, renal function (eGFR, UACR), LFTs, BP, weight/BMI/waist, retinal photography, foot exam, ECG, thyroid (if indicated), coeliac (if indicated), urinary albumin:creatinine ratio
    \item \textbf{C-peptide / autoantibodies:} if diagnostic uncertainty (lean, young, ketotic)
    \item \textbf{CVD risk:} Australian absolute CVD risk calculator (age, sex, smoking, BP, cholesterol, diabetes, HDL); consider coronary calcium score if intermediate
    \item \textbf{Complication screening (annual):} HbA1c 3--6 monthly, retinal photography, UACR/eGFR, lipids, BP, foot risk assessment, psychosocial (diabetes distress, depression)
    \item \textbf{NAFLD assessment:} FIB-4, ELF, FibroScan if LFTs elevated or metabolic risk factors
    \item \textbf{OSA screening:} STOP-BANG if symptoms
\end{itemize}

\section*{Management}
\begin{itemize}
    \item \textbf{Structured care:} GP Management Plan + Team Care Arrangement (Medicare); annual cycle of care; diabetes educator, dietitian, exercise physiologist, podiatrist, psychologist as needed
    \item \textbf{Lifestyle intervention (foundation, all patients):}
    \begin{itemize}
        \item Weight loss: 5--10\% improves glycaemia, BP, lipids; 15\%+ can induce remission (DiRECT trial)
        \item Diet: Mediterranean, low-carb, DASH, VLED (800 kcal, medically supervised); individualised; dietitian-led; carbohydrate awareness
        \item Physical activity: 150--300 min/week moderate aerobic + 2 resistance sessions; reduce sedentary time; exercise physiologist
        \item Sleep: 7--9 h; treat OSA
        \item Smoking cessation: critical
        \item Alcohol: $\le$10 std drinks/week, $\le$4/day
    \end{itemize}
    \item \textbf{Pharmacotherapy (individualised, per ADA/EASD/ADS guidelines):}
    \begin{itemize}
        \item \textbf{Metformin:} first-line unless contraindicated (eGFR <30, intolerance); weight neutral, low hypoglycaemia risk, cardiovascular safety, low cost; extended-release for GI intolerance
        \item \textbf{GLP-1 receptor agonists (semaglutide, dulaglutide, liraglutide, tirzepatide [dual GIP/GLP-1]):} high efficacy (HbA1c -1.0--2.0\%), weight loss (5--15\%), cardiovascular benefit (MACE reduction), renal benefit; injectable (weekly/daily); oral semaglutide available; PBS for T2D with HbA1c >7\% on metformin
        \item \textbf{SGLT2 inhibitors (empagliflozin, dapagliflozin, ertugliflozin):} HbA1c -0.5--0.8\%, weight loss 2--3 kg, CVD/HF/renal benefit (landmark trials: EMPA-REG, DECLARE, CREDENCE, DAPA-CKD, EMPEROR); genital mycotic infections, euglycaemic DKA risk (educate, hold if ill/surgery), avoid if eGFR <20 (dapagliflozin/empagliflozin approved down to 20 for HF/CKD)
        \item \textbf{DPP-4 inhibitors (sitagliptin, linagliptin, saxagliptin, alogliptin):} weight neutral, low hypoglycaemia risk, modest efficacy (-0.5--0.7\%), no CV harm (saxagliptin/alogliptin: HF caution); oral; linagliptin no renal dose adjustment
        \item \textbf{Sulfonylureas (gliclazide MR, glimepiride, glipizide):} inexpensive, effective (-1.0--1.5\%), but hypoglycaemia risk, weight gain, no CV benefit, $\beta$-cell exhaustion; second/third line; avoid in elderly/CKD
        \item \textbf{Thiazolidinediones (pioglitazone):} insulin sensitiser, durable, weight gain, oedema, fracture risk, bladder cancer concern, HF contraindicated; niche use
        \item \textbf{Insulin:} basal (glargine U100/U300, degludec, detemir) when oral agents insufficient; add prandial rapid-acting if needed; start low (10 U or 0.1--0.2 U/kg), titrate; CGM for insulin-treated; biosimilars available
        \item \textbf{Combination:} early dual/triple therapy if HbA1c >9\% or >1.5\% above target; fixed-ratio combinations (GLP-1 + basal insulin: iGlarLixi, iDegLira)
    \end{itemize}
    \item \textbf{Cardiorenal protection:} SGLT2i + GLP-1 RA + statin + ACEi/ARB + BP control -- four pillars
    \item \textbf{Remission:} defined as HbA1c <6.5\% off glucose-lowering medications for $\ge$3 months; achievable with >10--15\% weight loss (VLED, bariatric surgery, GLP-1 RA); more likely early (<6 yr duration), preserved $\beta$-cell function
    \item \textbf{Bariatric surgery:} BMI $\ge$35 with comorbidity or $\ge$30 with T2D; sleeve gastrectomy / RYGB; T2D remission 60--80\% at 2 years; durable with weight maintenance
    \item \textbf{Technology:} CGM (PBS for insulin-treated T2D); smart pens; apps for carb counting, data sharing
\end{itemize}

\section*{Complications}
\begin{itemize}
    \item \textbf{Microvascular:} retinopathy (NPDR $\to$ PDR, DME), nephropathy (albuminuria $\to$ CKD $\to$ ESRD), neuropathy (distal symmetric, autonomic, mononeuropathies)
    \item \textbf{Macrovascular:} CVD (MI, HF, stroke) -- 2--4x risk; PAD; T2D = coronary equivalent
    \item \textbf{Other:} NAFLD/NASH $\to$ cirrhosis/HCC; OSA; infections (fungal, bacterial, COVID severity); cognitive decline/dementia (1.5--2x); depression (2x); cancer (liver, pancreas, colorectal, breast, endometrial); hearing loss; frozen shoulder; foot ulcer/amputation
    \item \textbf{Acute:} hypoglycaemia (insulin/SU), HHS, DKA (SGLT2i, severe illness), lactic acidosis (metformin in severe renal/hepatic impairment)
    \item \textbf{Pregnancy:} pre-existing T2D: preconception HbA1c <6.5\%, folic acid 5 mg, review medications (stop SGLT2i, GLP-1 RA, statin, ACEi/ARB), specialist care
\end{itemize}

\section*{Prognosis and Outcomes}
T2D reduces life expectancy by ~6 years (diagnosis at 50) to ~10 years (diagnosis at 30). Cardiovascular disease is the leading cause of death. Intensive multifactorial intervention (Steno-2) reduces microvascular complications by 50\%, macrovascular by 50\%, mortality by 45\% over 20 years. Modern therapies (GLP-1 RA, SGLT2i) provide cardiorenal protection independent of glucose lowering. Remission is a realistic goal for many, especially with early intensive lifestyle or metabolic surgery. The treatment landscape has shifted from glucocentric to organ-protective.

\section*{Prevention and Self-Care}
\begin{itemize}
    \item \textbf{Primary prevention:} AUSDRISK screening; lifestyle intervention (Diabetes Prevention Program: 58\% reduction) for high-risk; metformin considered for BMI $\ge$35, age <60, GDM history
    \item \textbf{Weight management:} prevent gain from young adulthood; act early on small gains
    \item \textbf{Diet:} Mediterranean/low-carb/VLED; dietitian support
    \item \textbf{Activity:} move more, sit less; resistance training preserves muscle during weight loss
    \item \textbf{Monitoring:} HbA1c 3--6 monthly; home glucose if on insulin/SU; CGM if eligible
    \item \textbf{Medication adherence:} understand purpose; discuss side effects; don't stop without advice
    \item \textbf{Annual cycle of care:} eyes, kidneys, feet, heart, lipids, BP, teeth, mental health
    \item \textbf{Sick-day plan:} hold SGLT2i, metformin (if dehydrated), SU; continue insulin; monitor glucose/ketones; hydrate; seek care early
    \item \textbf{Foot care:} daily inspection; professional assessment; offloading if ulcer
    \item \textbf{Support:} Diabetes Australia, NDSS, Diabetes NSW\&ACT / Diabetes Victoria / Diabetes WA / Diabetes SA / Diabetes Queensland / Diabetes Tasmania, Australian Diabetes Educators Association
\end{itemize}

\section*{Sources and Further Reading}
The following reputable sources provide further information for healthcare students and patients seeking reliable, current advice about type 2 diabetes.

\begin{itemize}
    \item Diabetes Australia. \textit{National Diabetes Services Scheme (NDSS), resources, support.} https://www.diabetesaustralia.com.au/
    \item Australian Diabetes Society (ADS). \textit{Clinical guidelines, position statements.} https://www.ads.org.au/
    \item RACGP. \textit{Diabetes management in general practice (Diabetes Handbook).} https://www.racgp.org.au/
    \item American Diabetes Association (ADA) / European Association for the Study of Diabetes (EASD). \textit{Consensus reports on T2D management.} https://diabetesjournals.org/care/ / https://www.easd.org/
    \item National Institute for Health and Care Excellence (NICE). \textit{Type 2 diabetes guidelines.} https://www.nice.org.uk/
    \item Healthdirect Australia. \textit{Type 2 diabetes.} https://www.healthdirect.gov.au/type-2-diabetes
    \item Baker Heart and Diabetes Institute. \textit{Research, resources, clinical services.} https://www.baker.edu.au/
    \item Therapeutic Guidelines (Australia). \textit{Diabetes mellitus.} https://www.tg.org.au/
\end{itemize}